% SEGMENT OLP-0425-S01
% Part: normal-modal-logic
% Chapter: frame-correspondence
% Section: equivalence-S5

\documentclass[../../../include/open-logic-section]{subfiles}

\begin{document}

\olfileid{nml}{frd}{es5}

\olsection{சமானத் தொடர்புகளும் \Log{S5} உம்}

அனைத்துத் தொடர்பைக் கொண்ட எல்லாச் சட்டகங்களிலும், அதாவது ஒவ்வோர் உலகையும்
தன்னையும் சேர்த்து ஒவ்வோர் உலகிலிருந்தும் அணுக முடியும் சட்டகங்கள் அனைத்திலும்,
செல்லுபடியாகும் !!{formula}s இன் கணமாக வாய்ப்புநிலைத் தருக்கம்~\Log{S5}
சிறப்பிக்கப்படுகிறது. அத்தகைய சூழலில் $\Box$ இன்றியமையாமைக்கும் $\Diamond$
சாத்தியத்தன்மைக்கும் ஒத்திருக்கின்றன: $!A$ \emph{ஒவ்வோர்} உலகிலும் மெய்யாக
இருந்தால் $\Box !A$ மெய்; $!A$ \emph{ஏதோ} ஓர் உலகில் மெய்யாக இருந்தால்
$\Diamond !A$ மெய். தற்சுட்டு, சமச்சீர், கடப்புப் பண்புகளைக் கொண்ட எல்லாச்
சட்டகங்களிலும், அதாவது எல்லாச் \emph{சமானத் தொடர்புகளிலும்}, செல்லுபடியாகும்
!!{formula}s ஆகவும் \Log{S5} ஐச் சிறப்பிக்கலாம் என்று தெரியவருகிறது.

\begin{defn}
  $W$ மீதான ஈருறுப்பு உறவு~$R$ தற்சுட்டு, சமச்சீர், கடப்பு ஆகிய பண்புகளைப்
  பெற்றிருந்தால், எனில் மட்டுமே, அது ஒரு \emph{சமானத் தொடர்பு}. எல்லா
  $u,v \in W$ க்கும் $Ruv$ எனில், எனில் மட்டுமே, $W$ மீதான தொடர்பு~$R$
  \emph{அனைத்துத் தொடர்பு}.
\end{defn}

\Ax{T}, \Ax{B}, \Ax{4} ஆகியவை முறையே தற்சுட்டு, சமச்சீர், கடப்புச்
சட்டகங்களைச் சிறப்பிப்பதால், அணுகுறவு ஒரு சமானத் தொடர்பாக இருக்கும்
சட்டகங்களே இந்த மூன்று !!{formula}s உம் செல்லுபடியாகும் சட்டகங்கள்.
சமானத் தொடர்புகளை வரையறுக்க நாம் பயன்படுத்திய நிபந்தனைகள்,~$R$ மீதான மற்ற
அறிமுகமான நிபந்தனைகளின் சேர்க்கைகளுக்கு நிகராக இருப்பதால், வேறு
!!{formula}s இன் சேர்க்கைகளாலும் சமானத் தொடர்புகளைச் சிறப்பிக்கலாம் என்று
தெரியவருகிறது.

\begin{prop}\ollabel{prop:equivalences}
  பின்வருவன ஒன்றுக்கொன்று நிகரானவை:
  \begin{enumerate}
  \item $R$ ஒரு சமானத் தொடர்பு;
  \item $R$ தற்சுட்டு, யூக்ளிடியப் பண்புகளைக் கொண்டது;
  \item $R$ தொடர், சமச்சீர், யூக்ளிடியப் பண்புகளைக் கொண்டது;
  \item $R$ தொடர், சமச்சீர், கடப்புப் பண்புகளைக் கொண்டது.
  \end{enumerate}
\end{prop}

\begin{proof}
  பயிற்சி.
\end{proof}

\begin{prob}
  பின்வருவனவற்றைக் காட்டுவதன் மூலம்
  \olref[nml][frd][es5]{prop:equivalences} ஐ நிறுவுக:
  \begin{enumerate}
  \item $R$ சமச்சீர்ப் பண்பும் கடப்புப் பண்பும் கொண்டிருந்தால், அது
    யூக்ளிடியப் பண்புடையது.
  \item $R$ தற்சுட்டுப் பண்புடையதாக இருந்தால், அது தொடர் பண்புடையது.
  \item $R$ தற்சுட்டு, யூக்ளிடியப் பண்புகளைக் கொண்டிருந்தால், அது
    சமச்சீர்ப் பண்புடையது.
  \item $R$ சமச்சீர், யூக்ளிடியப் பண்புகளைக் கொண்டிருந்தால், அது
    கடப்புப் பண்புடையது.
  \item $R$ தொடர், சமச்சீர், கடப்புப் பண்புகளைக் கொண்டிருந்தால், அது
    தற்சுட்டுப் பண்புடையது.
  \end{enumerate}
  நிபந்தனைகள் நிகரானவை என்ற நிறுவலுக்கு இது ஏன் போதுமானது என்று விளக்குக.
\end{prob}

\olref{prop:equivalences} ஆனது \olref[prf][prs]{prop:S5} இன் பொருண்மையியல்
இணையாகும்; ஏனெனில் $R$ ஒரு சமானத் தொடர்பாக இருக்கும் சட்டகங்களின்
வாய்ப்புநிலைத் தருக்கத்திற்கு (மரபாக~\Log{S5} என்று குறிப்பிடப்படும்
தருக்கத்திற்கு) அது ஒரு நிகரான சிறப்பிப்பைக் கொடுக்கிறது.

அனைத்துத் தொடர்புகளுக்கும் சமானத் தொடர்புகளுக்கும் இடையிலான உறவு என்ன?
ஒவ்வோர் அனைத்துத் தொடர்பும் ஒரு சமானத் தொடர்பாக இருந்தாலும், ஒவ்வொரு
சமானத் தொடர்பும் அனைத்துத் தொடர்பல்ல என்பது தெளிவு. எனினும், எல்லா
அனைத்துத் தொடர்புகளிலும் செல்லுபடியாகும் !!{formula}s, எல்லாச் சமானத்
தொடர்புகளிலும் செல்லுபடியாகும் வாய்பாடுகளோடு துல்லியமாக ஒத்தவை.

\begin{prop}
  $R$ ஒரு சமானத் தொடர்பாக இருக்கட்டும். ஒவ்வொரு $w \in W$ க்கும்,
  $w$ இன் \emph{சமான வகுப்பை} $[w] = \{w'\in W : Rww'\}$ என்ற கணமாக
  வரையறுக்க. அப்போது:
  \begin{enumerate}
  \item $w \in [w]$;
  \item ஒவ்வொரு சமான வகுப்பு~$[w]$ மீதும் $R$ அனைத்துத் தொடர்பாகும்;
  \item சமான வகுப்புகளின் தொகுப்பு, $W$ ஐ ஒன்றையொன்று விலக்கும், கூட்டாக
    முழுமையாக்கும் உட்கணங்களாகப் பிரிக்கிறது.
  \end{enumerate}
\end{prop}

\begin{prop}\ollabel{prop:S5=univ}
  $R$ ஒரு சமானத் தொடர்பாக இருக்கும் எல்லாச் சட்டகங்கள் $\mModel{F}
  =\tuple{W,R}$ இலும் !!^a{formula}~$!A$ செல்லுபடியாக இருந்தால், எனில்
  மட்டுமே, $R$ அனைத்துத் தொடர்பாக இருக்கும் எல்லாச் சட்டகங்கள்
  $\mModel{F} =\tuple{W,R}$ இலும் அது செல்லுபடியாகும். ஆகவே அனைத்துத்
  தொடர்புச் சட்டகங்களின் தருக்கம்~\Log{S5} தான்.
\end{prop}

\begin{proof}
  $W$ மீதான அனைத்துத் தொடர்பு~$R$ ஒரு சமானத் தொடர்பு என்பதை உடனடியாகச்
  சரிபார்க்கலாம். ஆகவே $R$ ஒரு சமானத் தொடர்பாக இருக்கும் எல்லாச்
  சட்டகங்களிலும் $!A$ செல்லுபடியாக இருந்தால், அது அனைத்துத் தொடர்புச்
  சட்டகங்கள் அனைத்திலும் செல்லுபடியாகும். மற்ற திசையை முரண்மாற்றாக
  வாதிடுகிறோம்: $W$ மீதான சமானத் தொடர்பு~$R$ ஐக் கொண்ட சட்டகம்
  $\tuple{W,R}$ ஐ அடிப்படையாகக் கொண்ட மாதிரி $\mModel{M} = \tuple{W,
  R, V}$ இன் ஓர் உலகம்~$w$ இல் தவறும் ஒரு !!a{formula}~$!B$ ஐக் கொள்க.
  ஆகவே $\mSat/{M}{!B}[w]$. மாதிரி $\mModel{M}' = \tuple{W', R', V'}$ ஐப்
  பின்வருமாறு வரையறுக்க:
  \begin{enumerate}
  \item $W' = [w]$;
  \item $W'$ மீது $R'$ அனைத்துத் தொடர்பு;
  \item $V'(p) = V(p) \cap W'$.
  \end{enumerate}
  (எனவே $\mModel{M}'$ இலுள்ள உலகங்களின் கணம்~$W'$, \olref{fig:partition}
  இல் நிழலிட்ட பகுதியால் குறிப்பிடப்படுகிறது.) $W'$ மீது $R$, $R'$ ஆகியவை
  ஒத்துப்போகின்றன என்பதை எளிதாகக் காணலாம். பின்னர் !!{formula}s மீதான
  தொகுத்தறிதலால், ஒவ்வொரு~$!A$ க்கும் எல்லா $w' \in W'$ க்கும்,
  $\mSat{M'}{!A}[w']$ எனில், எனில் மட்டுமே, $\mSat{M}{!A}[w']$ என்று
  காட்டலாம் ($W' \subseteq W$ என்பதால் இது பொருளுடையது). குறிப்பாக,
  $\mSat/{M'}{!B}[w]$; மேலும், அனைத்துத் தொடர்புச் சட்டகத்தை அடிப்படையாகக்
  கொண்ட ஒரு மாதிரியில் $!B$ தவறுகிறது.
\end{proof}

\begin{figure}[t]
  \centering
  \begin{tikzpicture}[node distance=2cm, auto, thick]
    \clip [rounded corners] (0,0) -- (6,0) -- (6,4) -- (0,4) --  cycle;
    \begin{scope}
      \clip (0,0) -- (2,0)  .. controls (1.5,1.5) and (4.5,1.5)
      .. (4,4) -- (0,4) -- (0,0) -- cycle;
      \filldraw[fill=gray!40] (0,4) -- (0,2) .. controls (1.5,1.5)
      and (4.5,1.5) .. (4.5,0) -- (6,0) -- (6,4) -- (0,4) --  cycle;
    \end{scope}
    \draw [name path=line1] (2,0) .. controls (1.5,1.5) and (4.5,1.5) .. (4,4);
    \draw [name path=line2] (0,2)  .. controls (1.5,1.5) and (4.5,1.5) .. (4.5,0);
    \path [name intersections={of=line1 and line2}];
    \draw [rounded corners, use as bounding box, very thick] (0,0)
    -- (6,0) -- (6,4) -- (0,4) --  cycle;
    \node at (2,3) {$[w]$} ;
    \node at (1,1) {$[u]$} ;
    \node at (3,0.75) {$[v]$} ;
    \node at (4.75,2) {$[z]$} ;
  \end{tikzpicture}
  \caption{$W$ ஐச் சமான வகுப்புகளாகப் பிரித்தல்.}
  \ollabel{fig:partition}
\end{figure}

\end{document}
